% ch2_d §2.8–2.10 (书页 59–67 / p073–p081)
%--C-TAIL--
布里渊区作为倒格子的一个晶胞，可以重复铺排而恰好盖满矢量 $\mathbf{K}$ 的整个空间；$\mathbf{K}$ 的任何一个值都可以唯一地由一个 $\mathbf{g}$ 值——倒格子的一个平移矢——和一个 $\mathbf{q}$ 值构造出来，$\mathbf{q}$ 是一个被限制在该格子一个晶胞之内的矢量。用 §1.5 的语言来说，$\mathbf{q}$ 就是 $\mathbf{K}$ \emph{约化}（reduced）到第一布里渊区后的值。

因此，在这两种情形下式 (2.98) 都成立，其中 $\mathbf{q}$ 由 (2.99) 或 (2.100) 给出，视何者适用而定；如果我们把 $\mathbf{g}=0$ 也包括进倒格矢的完全集合之中，则关系式 (2.100) 恒成立。无论我们选取哪个方向，总会有一些散射，只不过其振幅将取决于碰巧具有满足这些条件的正确波矢的那些晶格振动的振幅。

\section{声子}

一个散射过程能否发生，并不单由矩阵元异于零所决定；还有能量条件需要满足。在理想晶格的情形，整个系统是静态的，因此初态和末态的电子（X 射线、中子）具有相同的能量。

的确，晶格处于不断的运动之中。(2.94) 中的矢量 $\mathbf{U}_{\mathbf{q}}$ 应当含有形如 $\exp(i\nu_{\mathbf{q}}t)$ 的时间因子，其中 $\nu_{\mathbf{q}}$ 是波矢为 $\mathbf{q}$ 的简正模的频率。实际上，$\mathbf{U}_{\mathbf{q}}$ 是三个这类矢量之和，它们的振幅、方向和频率各不相同，对应于该波矢的三个不同的声学模。但我们一次只取其中一个，并且记住，入射束与散射束的态 $\Psi_{\mathbf{k}}$ 和 $\Psi_{\mathbf{k}'}$ 必须具有能量 $\mathcal{E}(\mathbf{k})$ 和 $\mathcal{E}(\mathbf{k}')$，从而带有形如 $\exp\{i\mathcal{E}(\mathbf{k})t/\hbar\}$ 与 $\exp\{i\mathcal{E}(\mathbf{k}')t/\hbar\}$ 的时间因子。当我们对时间求平均时，将会遇到如下形式的积分
\begin{equation*}
\int \exp[i\{\mathcal{E}(\mathbf{k})-\mathcal{E}(\mathbf{k}')+\hbar\nu_{\mathbf{q}}\}t/\hbar]\,dt,
\end{equation*}
除非
\begin{equation*}
\mathcal{E}(\mathbf{k})-\mathcal{E}(\mathbf{k}')+\hbar\nu_{\mathbf{q}}=0.
\tag{2.101}
\end{equation*}
否则它们恒为零。换言之，衍射是\emph{非弹性}的；束流获得了与之相互作用的那个晶格模的一个能量量子。

其实这还不完全，因为晶格振动中还会出现共轭的时间依赖 $\exp(-i\nu_{\mathbf{q}}t)$，从而给出电子损失能量 $\hbar\nu_{\mathbf{q}}$ 的衍射。在这两种情形下，我们上面勾画的论证都可以借助于标准的时间依赖微扰理论而变得完全严格。

条件 (2.100) 与 (2.101) 可以给出一种非常直观的解释。我们说，入射电子（或 X 射线、中子）与晶格相互作用而消灭（或产生）了一个\emph{声子}（phonon），其波矢为 $\mathbf{q}$，能量为 $\hbar\nu_{\mathbf{q}}$。我们用图 33$(a)$ 或 $(b)$ 那样的图形来表示这些过程。条件 (2.101) 就是该过程中的能量守恒规则。

如果我们考察对第一区中 $(\mathbf{k}'-\mathbf{k})$ 的一切值都成立的条件 (2.99)，就会发现它看上去像一条动量守恒定律。乘以 $\hbar$，便得
\begin{equation*}
\hbar\mathbf{k}'=\hbar\mathbf{k}+\hbar\mathbf{q}.
\tag{2.102}
\end{equation*}
电子的动量吸收了被消灭的那个声子的\emph{晶体动量}（crystal momentum）$\hbar\mathbf{q}$。这正是为“声子”这个名称辩护的论证——声子是一种具有“类粒子”性质的量子化声激发，与“光子”相类比。

%FIG{33}: 电子散射过程：(a) 声子吸收；(b) 声子发射。

但在一般情形下，这条“动量守恒规则”并不成立；正如 (2.100) 所表明的，除声子的动量之外，电子还可以获得或失去动量 $\hbar\mathbf{g}$。我们称这样的过程为倒逆过程（Umklapp process），或 U 过程（U-process）。特殊情形 (2.102) 则可称为正常过程（Normal process），或 N 过程（N-process）。

这个结果其实并不令人惊讶。额外的动量 $\hbar\mathbf{g}$ 只不过是传给了作为整体的晶体。在构造晶格振动态时，我们忽略了理想晶格质心的运动——晶格振动正是相对于它来度量的。可以查验，只要把与这个自由度相联系的运动包括进来，基本的动量守恒定律总体上仍然得到满足。一个倒逆过程可以看成是在产生（或消灭）一个声子的同时发生了一次 Bragg 反射；后一过程中的动量显然传给了作为整体的晶体。

非弹性衍射现象是研究晶体晶格动力学的一种极有价值的工具。沿特定方向衍射的束流与具有确定波矢 $\mathbf{q}$ 的晶格模相联系。我们可以考察衍射粒子能量的变化，从而测出 $\hbar\nu_{\mathbf{q}}$。通过朝不同方向观察，并把晶体转到不同的取向，就可以把整个函数 $\nu_{\mathbf{q}}$ 测绘出来。当然，它有若干个不同偏振的分支，但通过系统分析可以把它们分开。中子对局部磁矩的敏感性（§2.6）意味着，对极化中子束被磁振子（magnons）非弹性衍射的情形也有一个类似的理论（见 §10.10），从而给出关于磁振子色散性质的类似信息。

然而，这个实验只有用“热”中子才切实可行；这种中子在能量约为 0.1 eV 时，其波长与晶格间距同数量级。由声子能量引起的频移——其数量级为 $k\Theta$ 或更小，也许 0.01 eV——很容易观察到。对于电子和 X 射线，束流的能量必须高得多——几十或几百电子伏——因此衍射过程中微小的能量变化无法被探测到。

晶体动量的选择定则 (2.102) 依赖于晶格的长程序。当这种有序不存在时，比如在液体或玻璃中，我们会观察到什么？弹性衍射将依赖于矩阵元 (2.86) 的平方，后者正比于结构因子的平方：
\begin{align*}
S(\mathbf{K}) &\equiv \left|\frac{1}{N}\sum_{\mathbf{l}} e^{-i\mathbf{K}\cdot\mathbf{R}_{\mathbf{l}}}\right|^{2}\\
&= \frac{1}{N^{2}}\sum_{\mathbf{l}\mathbf{l}'} e^{-i\mathbf{K}\cdot(\mathbf{R}_{\mathbf{l}}-\mathbf{R}_{\mathbf{l}'})}\\
&= \frac{1}{N}\int P(\mathbf{R})\,e^{-i\mathbf{K}\cdot\mathbf{R}}\,\mathrm{d}\mathbf{R}.
\tag{2.103}
\end{align*}
衍射图样依赖于对关联函数（pair correlation function）$P(\mathbf{R})$——即在相距 $\mathbf{R}$ 处找到两个原子的概率——的 Fourier 变换。在液体中，这正是我们熟悉的体系的径向分布函数（radial distribution function）。

对这个论证作一个简单的推广即可证明，非弹性衍射的振幅 $S(\mathbf{K},\nu)$ 必定是含时关联函数（time-dependent correlation function）$P(\mathbf{R},t)$ 的双重 Fourier 变换，$P(\mathbf{R},t)$ 度量的是这样的概率：已知在时刻 0 有一个原子（可能就是同一个）位于原点 0，在时刻 $t$ 于点 $\mathbf{R}$ 处找到一个原子的概率。例如，中子衍射可以测量波矢为 $\mathbf{K}$ 的密度涨落的频谱——这个物理量恰好只有在相当完善的晶格中才相当尖锐。类似地，极化中子会被自旋密度（spin density）涨落所衍射，这些涨落由类似的关联函数或它们在空间或动量、时间或能量上的 Fourier 变换所定义。

\section{德拜–沃勒因子}

(2.97) 中那些对应于含一个声子的散射过程的项，并不是结构因子 (2.95) 中可能出现的一切。容易看出，类似 (2.96) 的那些因子的乘积，以及指数展开中更高阶的项，会产生包含形如 $\mathbf{U}_{\mathbf{q}}\exp(i\mathbf{q}\cdot\mathbf{l})$ 的因子的各种乘积的贡献。对每一个这样的因子，我们都称相应的声子被产生或消灭了，因此这些项属于多声子过程（multiphonon processes）。

总的来说，随着阶数升高，这些过程的速率迅速下降，对非弹性衍射背景的贡献并不很大。然而，有一类重要的项来自展开式 (2.96) 中的平方项 $-|\mathbf{K}\cdot\mathbf{U}_{\mathbf{q}}|^{2}$，它们平均之后并不给出小的贡献。考察 (2.95) 与 (2.96) 即可发现，无论是弹性还是非弹性散射，矩阵元的平方都应当乘以
\begin{equation*}
e^{-2W}=\prod_{\mathbf{q}}\left\{1-|\mathbf{K}\cdot\mathbf{U}_{\mathbf{q}}|^{2}\right\},
\tag{2.104}
\end{equation*}
现在 $\mathbf{q}$ 则重新取遍整个布里渊区。

这个量称为 Debye–Waller 因子（Debye–Waller factor）。把它写成这种形式，是因为我们可以利用一条标准的代数定理
\begin{equation*}
\lim_{N\to\infty}\prod_{n=1}^{N}\left(1-\frac{1}{N}a_{n}\right)=\exp\left\{-\lim_{N\to\infty}\frac{1}{N}\sum_{n=1}^{N}a_{n}\right\},
\tag{2.105}
\end{equation*}
把乘积变换成求和。于是
\begin{equation*}
e^{-2W}=\exp\left\{-\sum_{\mathbf{q}}\frac{1}{2}\,|\mathbf{K}\cdot\mathbf{U}_{\mathbf{q}}|^{2}\right\},
\tag{2.106}
\end{equation*}
即
\begin{equation*}
W=\frac{1}{2}\sum_{\mathbf{q}}|\mathbf{K}\cdot\mathbf{U}_{\mathbf{q}}|^{2}.
\tag{2.107}
\end{equation*}
要计算这个和，我们需要知道第 $\mathbf{q}$ 个晶格模的振幅 $\mathbf{U}_{\mathbf{q}}$。它将是温度的函数。我们知道，该模中的平均能量由 (2.47) 给出：
\begin{equation*}
\bar{\mathcal{E}}_{\mathbf{q}}=\left(\bar{n}_{\mathbf{q}}+\frac{1}{2}\right)\hbar\nu_{\mathbf{q}},
\end{equation*}
其中 $\bar{n}_{\mathbf{q}}$ 是该模中声子的平均数，由 Bose–Einstein 公式 (2.46) 给出。

在经典处理中，我们可以把每个简谐振子模的能量算成其动能与势能之和。众所周知，这两者相等——于是由 (2.1) 和 (2.8) 得
\begin{align*}
\bar{\mathcal{E}}&=\sum_{s\mathbf{l}}M_{s}\,|\dot{\mathbf{u}}_{s\mathbf{l}}|^{2}\\
&=\sum_{s\mathbf{q}}NM_{s}\,|\dot{\mathbf{U}}_{s\mathbf{q}}|^{2}\\
&=\sum_{s\mathbf{q}}NM_{s}\,\nu_{\mathbf{q}}^{2}|\mathbf{U}_{s\mathbf{q}}|^{2}.
\tag{2.108}
\end{align*}
如果每个晶胞只有一个原子，质量为 $M$，则
\begin{align*}
|\mathbf{U}_{\mathbf{q}}|^{2}&=\bar{\mathcal{E}}_{\mathbf{q}}/NM\nu_{\mathbf{q}}^{2}\\
&=(\bar{n}_{\mathbf{q}}+\tfrac{1}{2})\hbar/NM\nu_{\mathbf{q}}.
\tag{2.109}
\end{align*}
我们知道 $\mathbf{U}_{\mathbf{q}}$ 在晶格谱每一支上的偏振，所以原则上能够精确计算出 Debye–Waller 因子。

为了看出它应当有怎样的行为，让我们假定一个三个模都具有相同速度的 Debye 模型。对任一偏振，平均而言我们应得
\begin{equation*}
|\mathbf{K}\cdot\mathbf{U}_{\mathbf{q}}|^{2}=\tfrac{1}{3}K^{2}|\mathbf{U}_{\mathbf{q}}|^{2},
\tag{2.110}
\end{equation*}
但若计入三个不同的偏振，因子 $\frac{1}{3}$ 便消去了。\footnote{有一条有用的规则：在这个模型中，三个模对每个 $\mathbf{q}$ 值都是简并的。因此可以随意选取偏振矢量，只要它们相互正交。于是可以取一个模的 $\mathbf{U}_{\mathbf{q}}$ 平行于 $\mathbf{K}$，另外两个垂直于 $\mathbf{K}$。这就给出了我们所需要的结果。}

由 (2.46)、(2.56) 和 (2.109) 得
\begin{align*}
W&=\frac{1}{2}K^{2}\frac{\hbar}{NM}\frac{1}{8\pi^{3}}\iiint\frac{\bar{n}_{\mathbf{q}}+\frac{1}{2}}{\nu_{\mathbf{q}}}\,d^{3}q\\
&=\frac{1}{2}\frac{\hbar^{2}K^{2}}{M}\int_{0}^{\nu_{D}}\left\{\frac{1}{e^{\hbar\nu/kT}-1}+\frac{1}{2}\right\}\frac{3\nu^{2}}{\hbar\nu\,\nu_{D}^{3}}\,d\nu\\
&=\frac{3}{2}\frac{\hbar^{2}K^{2}T^{2}}{Mk\Theta^{3}}\int_{0}^{\Theta/T}\left\{\frac{1}{e^{z}-1}+\frac{1}{2}\right\}z\,dz,\quad\text{正如式 (2.57) 中那样。}
\tag{2.111}
\end{align*}
高温下积分的上限很小，被积函数中的指数因子可以按 $z$ 的幂展开。结果是
\begin{equation*}
W\to\frac{3}{2}\frac{\hbar^{2}K^{2}T}{Mk\Theta^{2}}.
\tag{2.112}
\end{equation*}
于是，正比于矩阵元 $\mathcal{M}_{\mathbf{k}'\mathbf{k}}$ 的平方的 X 射线衍射图样，其强度将被如下因子所削减：
\begin{equation*}
e^{-2W}\sim\exp(-3\hbar^{2}K^{2}T/Mk\Theta^{2}).
\tag{2.113}
\end{equation*}
这个因子相当强地依赖于温度，也依赖于散射矢量的大小。如果我们对每个模的平均能量采用经典统计，即 $\bar{\mathcal{E}}_{\mathbf{q}}=kT$，也会得到这个结果。

在低于 Debye 温度的温度下，公式要更复杂一些，但我们注意到，在极低温下 $W$ 将趋于一个常数。这源于积分中的 $\frac{1}{2}$ 项——一项来自晶格的零点运动（zero-point motion）的贡献。这并不是一个可以忽略的效应：
\begin{equation*}
W\to\frac{3}{8}\frac{\hbar^{2}K^{2}}{Mk\Theta}\quad\text{当}\ T\to 0\ \text{时},
\tag{2.114}
\end{equation*}
它是 $W$ 在 $T\sim\Theta$ 处取值的 $\frac{1}{4}$。零点\emph{能}在物理上也许无足轻重，但与之相联系的\emph{运动}却可以被直接观察到。

非常类似的论证可以解释 Mössbauer 效应（Mössbauer effect）。原子核 γ 发射的谱中可能包含一条极窄的谱线，看不出它被原子在晶体中绕其平衡位置的热运动所展宽。这其实是一个简单的经典现象。以匀速 $v$ 运动的源所发射的频率为 $\omega$ 的辐射，会经受大小为 $\omega v/c$ 的 Doppler 频移。但设想这个源由一个束缚原子携带，该原子以频率 $\nu_{\mathbf{q}}$、振幅 $\mathbf{U}_{\mathbf{q}}$ 往复振动。由于速度现在随时间变化，我们观察到的发射辐射是调频的。正如调频（FM）广播的初等理论那样，谱现在分裂成位于 $\omega$、$\omega\pm\nu$、$\omega\pm2\nu$ 等频率处的若干条线。对于小的振动振幅，未发生改变的那条谱线在谱中所占的份额近似为 $1-\frac{1}{2}(\mathbf{k}\cdot\mathbf{U}_{\mathbf{q}})^{2}$，其中 $\mathbf{k}$ 是 γ 光子的波矢。像 (2.104–7) 那样把所有可能振动模的贡献合并起来，我们发现这条谱线的相对强度是一个有限量，由一个正像 Debye–Waller 因子一样的公式给出。

用固体理论的标准语言，我们可以说，γ 射线的发射传给原子的冲量有一部分被声子的产生所吸收，声子的能量则表现为谱线的展宽。但有有限比例的动量传给了作为整体的晶体，不伴随声子的产生（或吸收），因而只有无穷小的能量损失（或获得）。正如我们对 X 射线衍射的分析那样，Debye–Waller 因子度量的正是这类弹性即无反冲（recoilless）事件所占的比例。这条极锐的谱线当然极其有用：可以用来探测分子与固体中的磁场，检验相对论理论，等等。

(2.111) 的推导还顺带给了我们另一个有趣的结果。由 (2.1) 和 (2.8) 显然可见，在 Debye 温度以上，每个原子绕其格点振动的均方振幅由下式给出：
\begin{align*}
\frac{1}{N}\sum_{\mathbf{l}}|\mathbf{u}_{\mathbf{l}}|^{2}&=\sum_{\mathbf{q}}|\mathbf{U}_{\mathbf{q}}|^{2}\\
&\approx\frac{9\hbar^{2}T}{Mk\Theta^{2}}
\tag{2.115}
\end{align*}
于是在温度 $T$ 下，每个原子偏离其平衡位置的均方根位移将是，比如说，晶胞平均半径 $r_{s}$ 的一个分数 $x$，其中
\begin{equation*}
x=\sqrt{\frac{9\hbar^{2}T}{Mk\Theta^{2}r_{s}^{2}}}.
\tag{2.116}
\end{equation*}
Lindemann 熔化公式（Lindemann melting formula）正是基于这个观念。它认为，当 $x$ 达到某个标准值 $x_{m}$ 时，固体就必须熔化。于是，熔化温度 $T_{m}$ 便与固体的其他原子常数通过下式相联系：
\begin{equation*}
T_{m}=\frac{x_{m}^{2}}{9\hbar^{2}}\,Mk\Theta^{2}r_{s}^{2}.
\tag{2.117}
\end{equation*}
看来在大多数固体中 $x_{m}$ 处于 0.2–0.25 的范围内。

这条规则可以用来在已知 $T_{m}$ 的情形下近似地估计 Debye 温度。它也为估计像 $W$ 这类依赖于原子振动振幅的量的数值提供了一条便捷的捷径。例如，我们可以证明，(2.112) 可以写成
\begin{equation*}
W\approx x_{m}^{2}\,\frac{T}{T_{m}}\,\frac{K^{2}}{q_{D}^{2}},
\tag{2.118}
\end{equation*}
其中 $q_{D}$ 是 Debye 波数 (2.53)。

\section{非谐性与热膨胀}

在本章开始时，我们把晶体的势能展开成了晶格位移的 Taylor 级数。但这个级数——式 (2.2)——在二阶项处被截断了。还会有更进一步的项，例如
\begin{equation*}
\mathcal{V}^{(3)}=\frac{1}{3!}\sum_{\substack{ss's''\\ ll'l''\\ jj'j''}}u_{s\mathbf{l}}^{j}\,u_{s'\mathbf{l}'}^{j'}\,u_{s''\mathbf{l}''}^{j''}\left[\frac{\partial^{3}\mathcal{V}}{\partial u_{s\mathbf{l}}^{j}\,\partial u_{s'\mathbf{l}'}^{j'}\,\partial u_{s''\mathbf{l}''}^{j''}}\right]_{0},
\tag{2.119}
\end{equation*}
如此等等。这些系数的实际计算是一个非常复杂的问题，因为它们既涉及几何因子，又涉及原子间势的三阶导数。

有若干重要的物理现象与非谐项（anharmonic terms）相联系。其中最熟知的是热膨胀（thermal expansion）。要从 (2.119) 这类表达式直接导出它并不容易，但一般的物理图像却相当简单。随着温度升高，晶格振动的振幅增大，因而位移 $\mathbf{u}_{s\mathbf{l}}$ 等的均方根平均值增大。非谐项对晶体的自由能有贡献，于是自由能不再必然在围绕所假定的“平衡”位形（其中每个 $\mathbf{u}_{s\mathbf{l}}$ 均为零）的振动下取极小。整个晶体于是会膨胀（或收缩），直至找到总自由能取极小的那个体积。

由于缺乏关于非谐项的充分知识而无法作“第一性原理”的计算，我们可以唯象地假定晶格诸模的频率是体积的函数。为简单起见，设体积变化 $\Delta V$ 使每个模的频率产生相同的相对变化：
\begin{equation*}
\frac{\Delta\nu}{\nu}=-\gamma\,\frac{\Delta V}{V}.
\tag{2.120}
\end{equation*}
于是，晶体总自由能作为体积的函数便可写成
\begin{equation*}
F=\frac{1}{2}\,\frac{1}{\kappa}\left(\frac{\Delta V}{V}\right)^{2}+kT\sum_{\mathbf{q}}\ln\left(2\sinh\frac{\hbar\nu_{\mathbf{q}}}{2kT}\right).
\tag{2.121}
\end{equation*}
第一项是与作为弹性连续介质的固体的压缩率（compressibility）$\kappa$ 相联系的势能。第二项是诸晶格模自由能之和，由 Bose–Einstein 振子的常规统计力学给出。

对体积求微商，并利用 (2.120)，我们得到自由能取极小的条件：
\begin{align*}
\frac{1}{\kappa}\left(\frac{\Delta V}{V}\right)&=\sum_{\mathbf{q}}\gamma\hbar\nu_{\mathbf{q}}\,\tfrac{1}{2}\coth\frac{\hbar\nu_{\mathbf{q}}}{2kT}\\
&=\gamma\bar{\mathcal{E}}(T),
\tag{2.122}
\end{align*}
其中 $\bar{\mathcal{E}}$ 是温度 $T$ 下诸晶格模中的能量，如 (2.48) 所示。这样我们便得到了 Grüneisen 公式（Grüneisen formula）：温度 $T$ 下的体膨胀（dilatation）正比于平均热能量密度，即
\begin{equation*}
\frac{\Delta V}{V}=\kappa\gamma\bar{\mathcal{E}}(T).
\tag{2.123}
\end{equation*}
热膨胀系数（thermal expansion coefficient）作为体膨胀对温度的导数，正比于比热 $C_{V}$。

把这个公式与实验比较，可以定出 Grüneisen 常量（Grüneisen constant）$\gamma$，其值通常约为 2。它是刻画非谐性效应的一个方便的无量纲参量。在 Debye 模型中，我们可以把 (2.120) 写成
\begin{equation*}
\gamma=-\frac{\partial\ln\Theta}{\partial\ln V}
\tag{2.124}
\end{equation*}
的形式，它表明了体积对 Debye 温度的影响。说实话，作为热膨胀的一个模型，这实在过于简化了。体膨胀对不同模的影响方式并不相同；纵模的 $\gamma$ 值通常远大于横模的 $\gamma$ 值，因此在 (2.122) 中必须把它们作为单独的贡献分别计入。

非谐性的另一个效应表现在弹性常数上，弹性常数随体积和温度而变化。这些是复杂的现象，没有初等理论可循；它们依赖于若干个不同的参量。
